\newenvironment{novspacecenter}
{\parskip=0pt\par\nopagebreak\centering}
{\par\noindent\ignorespacesafterend}

\definecolor{charitegray}{RGB}{160 165 169}
\definecolor{charitegrayarea}{RGB}{229 229 229}

 \begin{titlepage}
	\begin{novspacecenter}
		\includegraphics[width=80mm]{Images/logo_inria.png}\hfill
		\includegraphics[width=80mm]{Images/logo_universite.png}
	\end{novspacecenter}
	\vskip19mm
	\begin{tikzpicture}[remember picture, overlay]         
	\fill [charitegrayarea, opacity=1] (-3.0,0.0) rectangle (19,-25);
	\end{tikzpicture}
	\raggedright
	\vskip18mm
	{
		\fontsize{22}{26.4}
		\bfseries
		\boldmath
%		\sffamily
		La méthode \textit{Shifted Boundary} pour un problème d'ondes linéaire avec zones de congestion
		\par
	}
	\vskip18mm
	{
		\fontsize{20}{24}
		\boldmath
		\textbf{Rapport de stage}
		\par
	}
	\vskip18mm
	{%
		\fontsize{15}{18}
		\boldmath
%		\sffamily
		\textsc{Pannetier Valentin}\\
		sous la direction de \textsc{Mr Mario Ricchiuto}.
		\par
	}
	\vfill
	{
		\large
		\vfill
			
		Inria Bordeaux-Sud-Ouest \& Université de Bordeaux - \today\\
		Année 2019-2020, Master Modélisation EDP
	}
\end{titlepage}

%\newpage